\chapter{Looking Beyond the Score Reveals Hidden Nuances}
\label{chap:beyond-score}

In Chapter~\ref{chap:right-metrics}, we saw the importance of defining metrics
that capture the properties we actually care about while controlling for confounding biases.
In this chapter, we focus
on the importance of looking beyond the score and analyzing the underlying success and failure 
modes of models. We present four examples showing how deeper analysis can reveal hidden nuances 
in the behavior of language models and advance our scientific understanding.

This chapter first discusses insights and analyses derived from three of the benchmarks 
introduced in the previous chapter. In Section~\ref{sec:cruxeval-beyond-score}, we dive deeper into 
the results from CRUXEval. For example, by comparing scores on CRUXEval with code generation benchmarks like HumanEval,
we discover that a cluster of distilled models (e.g. WizardCoder) perform much better on HumanEval compared
to CRUXEval than their base models, suggesting distillation may not fully transfer general programming abilities.
In Section~\ref{sec:counterfeit-beyond-score}, we analyze whether counterfeit samples have different 
properties depending on the difficulty of the problem or the strength of the model generating the samples. 
In Section~\ref{sec:livecodebench-beyond-score}, we return to the LiveCodeBench paper to analyze
overfitting to HumanEval, the role of post-training, and the gap between open and closed models.
Finally, in Section~\ref{sec:linc-beyond-score}, we introduce LINC, a neuro-symbolic approach to performing
logical reasoning. By analyzing the behavior of LINC compared to traditional chain-of-thought,
we discover that LINC and chain-of-thought have distinct failure modes, suggesting opportunities to
synthesize the two approaches. 


\section{CRUXEval: Generalization, Chain of Thought, Fine-Tuning, and Qualitative Failures}
\label{sec:cruxeval-beyond-score}
\subsection{Quantitative Analysis} \label{sec:quantitative-analysis}

\textbf{Correlation between scores on HumanEval and \benchmark}: After the release of Code Llama's model and GPT-3.5 and GPT-4's APIs, there have been many creative efforts to take data distilled from GPT models and use them to train more powerful code models such as WizardCoder \citep{luo2023wizardcoder}, Phi-1 \citep{gunasekar2023textbooks}, Phi-1.5 \citep{gunasekar2023textbooks}, and Phind \citep{Phind}. For example, WizardCoder 34B started with the Code Llama 34B base model and improved the HumanEval pass@1 score from 53.7\% to 73.2\%, a significant and impressive achievement. There remains curiosity about whether these models show more general improvements in other aspects of programming or code understanding \citep{gudibande2023false}. We measure this through \benchmark.

In Fig. \ref{fig:quant-analysis-humaneval-correlation-1}, we plot reported HumanEval scores (we did not reproduce them ourselves) against scores on \benchmark. Indeed, we spot some interesting outliers: when comparing the distilled models WizardCoder 34B and Phind 34B to \codellamalarge, we see that the distilled models score over 20\% more than Code Llama on HumanEval but do not show this drastic improvement when evaluated on both input and output predictions. In addition, the Phi-1 model outperforms most of the bigger models on HumanEval, but performs among the worst of all our evaluated models on \benchmark. Overall, this suggests that models optimized for the HumanEval task by distilling data from GPT-3.5 and GPT-4 (WizardCoder, Phind, Phi) may not have learned other code reasoning capabilities along the way. On the other hand, for models such as StarCoder, Mistral, CodeLlama, and DeepSeek-Base, we still see a positive trend between HumanEval score and \benchmark score, suggesting that code generation and execution/understanding abilities are correlated.

\begin{figure}[H]
     \centering
     \begin{subfigure}[b]{0.49\textwidth}
         \centering
         \includegraphics[scale=0.49]{figs/scatter/HumanEval_1_vs_input_pass1.pdf}
         \subcaption{\benchmarki vs. HumanEval}
         \label{fig:quant-analysis-humaneval-correlation-input}
     \end{subfigure}
     \hfill
     \begin{subfigure}[b]{0.49\textwidth}
         \centering
         \includegraphics[scale=0.49]{figs/scatter/HumanEval_1_vs_output_pass1.pdf}
         \subcaption{\benchmarko vs. HumanEval}
         \label{fig:quant-analysis-humaneval-correlation-output}
     \end{subfigure}
     \caption{Correlation between HumanEval pass@1 scores and \benchmarko pass@1 scores}
     \label{fig:quant-analysis-humaneval-correlation-1}
\end{figure}

\begin{tcolorbox}[colback=lightgreen, boxrule=0pt, arc=10pt, outer arc=10pt]
Base models show a weak correlation between HumanEval and \benchmark. For HumanEval, distilled models (WizardCoder, Phind, Phi) significantly beat their base models, but for \benchmark, no distilled model performs significantly better than \codellamalarge.
\end{tcolorbox}

\textbf{Relationship between input prediction and output prediction:}
In Fig. \ref{fig:quant-analysis-input-output-direct}, we compare the input prediction and output prediction pass@1 scores with each other. Conceptually, the two tasks seem relatively different: output prediction is directly testing code execution ability, while input prediction requires a higher-level understanding of the code's functionality. However, we discover that there is a strong correlation between their performance. This suggests the hypothesis that performance on relatively distinct coding-related tasks may be closely correlated. In addition, we see a relatively clear impact of scaling the model size on our two tasks.


\begin{figure}[H]
     \centering
     \begin{subfigure}[b]{0.49\textwidth}
         \centering
         \includegraphics[scale=0.49]{figs/scatter/output_1_vs_input_1.pdf}
         \subcaption{Input vs. Output Prediction, Direct}
         \label{fig:quant-analysis-input-output-direct}
     \end{subfigure}
     \hfill
     \begin{subfigure}[b]{0.49\textwidth}
         \centering
         \includegraphics[scale=0.5]{figs/scatter/cot_input_pass1_output_pass1.pdf}
         \subcaption{Output vs. Input Prediction, CoT}
         \label{fig:quant-analysis-input-output-cot}
     \end{subfigure}
     \caption{Correlation between Input and Output Prediction Scores, with and without CoT}
     \label{fig:quant-analysis-input-output}
\end{figure}

\begin{tcolorbox}[colback=lightgreen, boxrule=0pt, arc=10pt, outer arc=10pt]
With the potential exception of GPT models, performance on \benchmarki and \benchmarko seem to be very correlated. As the tasks seem relatively different, this suggests that the code reasoning capabilities of models may generalize from task to task.
\end{tcolorbox}


\textbf{Confusion matrix/error correlation for different models.}
Fig.~\ref{fig:error_corr-input-output} shows the pairwise correlation of pass@1 scores for each pair of models. The correlation is chosen based on its highest signal among cosine distance, Spearman and Kendall. The middle section of ``open''-ish models (StarCoder, Code Llama, DeepSeek, etc.) are strongly correlated with each other. Strong correlations are seen between sizes of the same model, between models of the same size, and between instruct and base (\textsc{Phind 34B, Wizard 34B} vs. \codellamalarge). CoT results also tend to have strong correlations with other CoT results, even GPT-4 vs Llama 13B. For the output task, \textsc{Deepseek} forms a small sub-cluster of especially strong associations.


\begin{figure}[H]
     \centering
     \begin{subfigure}[b]{0.49\textwidth}
         \centering
         \includegraphics[width=1\textwidth]{figs/error_correlation_input_heatmap.pdf}
     \end{subfigure}
     \hfill
     \begin{subfigure}[b]{0.49\textwidth}
         \centering
         \includegraphics[width=1\textwidth]{figs/error_correlation_output_heatmap.pdf}
     \end{subfigure}
     \caption{Correlation between predictions on input (left) and output (right)}
     \label{fig:error_corr-input-output}
\end{figure}


\begin{tcolorbox}[colback=lightgreen, boxrule=0pt, arc=10pt, outer arc=10pt]
Looking at model predictions, strong correlations are seen between sizes of the same model, between models of the same size, and between instruct and base models. Although what is hard for a better model tend to be hard for worse models on average, worse models succeeded on some examples where the better models fail completely.
\end{tcolorbox}

\subsection{Chain of Thought Prompting} \label{subsec:cot}
Next, we evaluate how the popular chain-of-thought (CoT) prompting method \citep{wei2022chain} affects the performance of Code Llama, GPT-3.5, and GPT-4 models on \benchmark. All results are reported using $N=10$ samples other than CodeLlama 13B and 34B without CoT, which are reported with $N=100$ samples. As before, pass@1 is reported with $T=0.2$ and pass@5 with $T=0.8$.

\textbf{Impact of CoT}: We begin by focusing our attention on the pass@1 scores of models with and without CoT. In Fig. \ref{fig:quant-analysis-input-output-cot}, we plot the input and output prediction scores of each model with and without CoT. First, GPT-4 benefits significantly more than other models. Second, output prediction boosts are generally larger than input prediction. In fact, CoT does not seem to improve Code Llama 13B and GPT-3.5 performance on input prediction. This is intuitive, as input prediction involves a more difficult reasoning task, while output prediction only requires executing the program step by step.

\begin{tcolorbox}[colback=lightgreen, boxrule=0pt, arc=10pt, outer arc=10pt]
CoT helps Code Llama 34B and GPT-4 on both input and output prediction, GPT-3.5 on only output prediction, and Code Llama 13B on neither task. CoT also leads to larger boosts on output prediction than input prediction. GPT-4 benefits significantly more from CoT than other models, achieving the highest pass@1 of 74.8\% on input prediction and 81.9\% on output prediction but still far from acing the benchmark.
\end{tcolorbox}

\textbf{CoT widens the gap between pass@5 and pass@1 scores}: In Fig. \ref{fig:pass5-vs-pass1-cot}, we plot the pass@5 scores against the pass@1 scores for all models. For models without CoT (shown in blue), there is a positive correlation between pass@1 and pass@5 scores. For models with CoT (shown in orange), we see an increase in the gap between pass@5 and pass@1 scores. This suggests that CoT increases generation diversity, which can improve the best-of-several score even when the first sample remains unreliable.

\begin{tcolorbox}[colback=lightgreen, boxrule=0pt, arc=10pt, outer arc=10pt]
Because CoT increases the diversity of generated inputs and outputs, models with CoT see a larger gap between pass@1 and pass@5 score compared to models without.
\end{tcolorbox}

\begin{figure}[H]
     \centering
     \begin{subfigure}[b]{0.49\textwidth}
         \centering
         \includegraphics[scale=0.45]{figs/scatter/input_pass_5v1.pdf}
         \subcaption{Input prediction}
         \label{fig:pass5-vs-pass1-cot-input}
     \end{subfigure}
     \hfill
     \begin{subfigure}[b]{0.49\textwidth}
         \centering
         \includegraphics[scale=0.45]{figs/scatter/output_pass_5v1.pdf}
         \subcaption{Output prediction}
         \label{fig:pass5-vs-pass1-cot-output}
     \end{subfigure}
     \caption{pass@5 score vs. pass@1 score with and without CoT}
     \label{fig:pass5-vs-pass1-cot}
\end{figure}

\textbf{Predictions of CoT vs. Base Model}: In Fig. \ref{fig:confusion-cot-all}, we show a confusion matrix over samples to better understand the correlations between direct output predictions and CoT predictions. For CodeLlama 13B, 34B, and GPT-3.5, we observe a large number of samples where direct prediction succeeds but CoT fails. However, with GPT-4, we observe that there are relatively few samples where this is the case.

\begin{figure}[H]
     \centering
     \begin{subfigure}[t]{0.49\textwidth}
         \centering
         \includegraphics[width=\textwidth]{figs/confusion_cot/confusion_codellama_13B_codellama_cot13B.pdf}
         \subcaption{Code Llama 13B}
         \label{fig:confusion-cot-codellama-13b}
     \end{subfigure}%
     \hfill
     \begin{subfigure}[t]{0.49\textwidth}
         \centering
         \includegraphics[width=\textwidth]{figs/confusion_cot/confusion_codellama_30B_codellama_cot30B.pdf}
         \subcaption{Code Llama 34B}
         \label{fig:confusion-cot-codellama-34b}
     \end{subfigure}
     \newline
     \newline
     \newline
     \begin{subfigure}[t]{0.49\textwidth}
         \centering
         \includegraphics[width=\textwidth]{figs/confusion_cot/confusion_gpt35_gpt35_cot.pdf}
         \subcaption{GPT-3.5}
         \label{fig:confusion-cot-gpt35}
     \end{subfigure}%
     \hfill
     \begin{subfigure}[t]{0.49\textwidth}
         \centering
         \includegraphics[width=\textwidth]{figs/confusion_cot/confusion_gpt4_gpt4_cot.pdf}
         \subcaption{GPT-4}
         \label{fig:confusion-cot-gpt4}
     \end{subfigure}
     \caption{Confusion Matrix of Direct Prediction vs. CoT Prediction $(T=0.2)$}
     \label{fig:confusion-cot-all}
\end{figure}


\begin{tcolorbox}[colback=lightgreen, boxrule=0pt, arc=10pt, outer arc=10pt]
While CoT generally improves performance overall, there are many individual samples where CoT actually hurts the prediction accuracy for Code Llama 13B/34B and GPT-3.5 on both input and output prediction. For GPT-4, CoT generally improves individual sample accuracy, more so for output prediction than for input prediction.
\end{tcolorbox}

\subsection{Fine-tuning Experiments} \label{subsec:finetuning}
Next, we do a preliminary analysis to understand the effect of simple fine-tuning schemes on \benchmark performance. We fine-tuned \codellamalarge on nearly 140K samples of Python functions distilled with the procedure outlined in Sec. \ref{sec:benchmark-construction}, without filtering. We perform weak decontamination, only removing samples where both the function and input-output pairs match samples in the benchmark. 

In particular, we finetune on a mixture of 50\% samples where the function is not in the benchmark and 50\% samples where the function is in the benchmark but input-output pairs are not, a very liberal setup. The training and testing accuracy over time is shown in Fig. \ref{fig:finetuning-accuracy-samples-plot-main}. Despite finetuning on programs very similar to the benchmark, we still observe a plateauing effect in the test accuracy, suggesting that our execution tasks may be too difficult to learn from this simple fine-tuning scheme.

\begin{figure}[H]
     \centering
     \begin{subfigure}[b]{0.48\textwidth}
         \centering
         \includegraphics[scale=0.4]{figs/finetuning/finetuning_input_moredata_accuracy.pdf}
         \subcaption{Input prediction}
         \label{fig:finetuning-accuracy-samples-plot-input-main}
     \end{subfigure}
     \hfill
     \begin{subfigure}[b]{0.48\textwidth}
         \centering
         \includegraphics[scale=0.4]{figs/finetuning/finetuning_output_moredata_accuracy.pdf}
         \subcaption{Output prediction}
         \label{fig:finetuning-accuracy-samples-plot-output-main}
     \end{subfigure}
     \caption{Improvements and Limits of \benchmark Performance after Fine-Tuning}
     \label{fig:finetuning-accuracy-samples-plot-main}
\end{figure}

\begin{tcolorbox}[colback=lightgreen, boxrule=0pt, arc=10pt, outer arc=10pt]
After fine-tuning on samples very similar to those in our benchmark, Code Llama 34B can match the performance of GPT-4 on both input and output prediction. However, accuracy plateaus at under 70\% for both tasks, so simple finetuning is far from solving the benchmark.
\end{tcolorbox}

\subsection{Qualitative Analysis} \label{sec:qualitative-analysis}
All models except GPT4 has over 50\% failure rate, showing they cannot do simple executions. 
In this section, we focus on GPT4 with CoT and verify that the remaining 20\% failures are due to the model, are consistent, and are indeed on \emph{simple} programs.

\paragraph{Failures of GPT-4 CoT.}
GPT-4 Cot scored 0/10 on 54 output prediction tasks and 65 input prediction tasks. On 22 problem, it scored 0/10 on both input and output prediction tasks. 
We manually check the 22 problems if they pass our criteria of being simple problems. Most are indeed simple (Listings~\ref{lst:gpt-4-simple1}, ~\ref{lst:gpt-4-simple2}). There are 2 problems that require counting to around 30 (Listing~\ref{lst:gpt-4-hard1}) and 2 problems (Listing~\ref{lst:gpt-4-hard2}) that require simulating a few actual steps, which might be difficult for direct generation but within scope for CoT. 

\begingroup
\lstset{basicstyle=\ttfamily\scriptsize,breakindent=0pt,postbreak={},columns=fullflexible,keepspaces=true}

\noindent\begin{minipage}[t]{.47\textwidth} % sample_218 steps: 37
\vspace{0pt}
\begin{lstlisting}[caption=GPT-4 has the right idea but cannot do the string concatenation correctly,label={lst:gpt-4-simple1}, breaklines=true, language=python]
def f(string, sep):
    cnt = string.count(sep)
    return((string+sep) * cnt)[::-1]
assert f('caabcfcabfc', 'ab') == 'bacfbacfcbaacbacfbacfcbaac'
# GPT4+CoT preds:
# f('baa', 'c')
# f('cba', 'f')
# 'bacfcabcfcaabcabfcabcfcaac'
# 'bacfbacfcbaabcabfbacfcbaabc'
\end{lstlisting}
\end{minipage}\hfill
\begin{minipage}[t]{.47\textwidth} % sample_209	steps: 26	
\vspace{0pt}
\begin{lstlisting}[caption=GPT-4 might have been misled by the variable name prefix,label={lst:gpt-4-simple2}, breaklines=true, language=python]
def f(prefix, s):
    return str.removeprefix(prefix, s)
assert f('hymi', 'hymifulhxhzpnyihyf') == 'hymi'
# GPT4+CoT preds:
# input
# f('p', 'phymi')
# f('', 'hymi')
# output
# 'fulhxhzpnyihyf'
# 'fulhxhzpnyihyf'
\end{lstlisting}
\end{minipage}


\noindent\begin{minipage}[t]{.47\textwidth} % sample_129	steps: 132	
\vspace{0pt}
\begin{lstlisting}[caption={GPT-4 CoT failures where solutions requires counting to 30}, label={lst:gpt-4-hard1}, breaklines=true, language=python]
def f(text, search_string):
    indexes = []
    while search_string in text:
        indexes.append(text.rindex(search_string))
        text = text[:text.rindex(search_string)]
    return indexes
assert f('ONBPICJOHRHDJOSNCPNJ9ONTHBQCJ', 'J') == [28, 19, 12, 6]

# GPT4+CoT preds:
# input
# f("0000123000012300001230000123", "123")
# f('bbbbbbabbbbbbabbbbbbbabbbbbbab', 'a')
# f("abcxxxxxxabcxxxxxxabcxxxxxxabc","abc")
# output
# [23, 13, 8, 5]
# [25, 18, 15, 11, 6]
# [7, 10, 14, 18]
\end{lstlisting}
\end{minipage}\hfill %sample_185	steps: 275	
\begin{minipage}[t]{.47\textwidth}
\vspace{0pt}
\begin{lstlisting}[caption={GPT-4 CoT failure, cannot easily tell the answer without running the loops}, label={lst:gpt-4-hard2}, breaklines=true, language=python]
def f(L):
    N = len(L)
    for k in range(1, N//2 + 1):
        i = k - 1
        j = N - k
        while i < j:
            # swap elements:
            L[i], L[j] = L[j], L[i]
            # update i, j:
            i += 1
            j -= 1
    return L
assert f([16, 14, 12, 7, 9, 11]) == [11, 14, 7, 12, 9, 16]
# GPT4+CoT preds:
# f([16, 9, 12, 7, 14, 11])
# f([16, 9, 7, 12, 14, 11])
# [11, 9, 7, 12, 14, 16]
# [11, 9, 7, 12, 14, 16]
\end{lstlisting}
\end{minipage}


Listings \ref{lst:gpt-4-output-fail1} and \ref{lst:gpt-4-output-fail2} show GPT-4 CoT failures for output prediction only. In Listing \ref{lst:gpt-4-output-fail2}, the model fails because it concludes that 6173 is not less than 1000. Small perturbations like changing to \texttt{0 < num and num < 1000} or changing the strings also failed. Both problems only have 2 possible answers and other models sometimes get them correctly whereas GPT4 CoT is consistently wrong. We manually tested scratchpad \cite{nye2021show} style prompts, which failed in the same way as regular CoT.

\noindent\begin{minipage}[t]{.47\textwidth}
\vspace{0pt}
\begin{lstlisting}[caption=GPT-4 CoT output,label={lst:gpt-4-output-fail1}, breaklines=true, language=python]
def f(text, suffix):
    if suffix == '':
        suffix = None
    return text.endswith(suffix)
assert f('uMeGndkGh', 'kG') == ??
# GPT-4 CoT: True
# should be False
\end{lstlisting}
\end{minipage}\hfill
\begin{minipage}[t]{.47\textwidth}
\vspace{0pt}
\begin{lstlisting}[caption=GPT-4 CoT output,label={lst:gpt-4-output-fail2}, breaklines=true, language=python]
def f(num):
    if 0 < num < 1000 and num != 6174:
        return 'Half Life'
    return 'Not found'
assert f(6173) == ??
# GPT-4 CoT: 'Half Life'
# should be 'Not found'
\end{lstlisting}
\end{minipage}

\textbf{Failures of GPT-4, Input Prediction}: Here are two simple failures on input prediction. Listings \ref{lst:gpt-4-input-fail1} and \ref{lst:gpt-4-input-fail2} show input prediction failures for concise and simple Python programs with and without CoT, respectively.

\noindent\begin{minipage}[t]{.47\textwidth}
\vspace{0pt}
\begin{lstlisting}[caption=GPT-4 CoT input,label={lst:gpt-4-input-fail1}, breaklines=true, language=python]
def f(text, repl):
    trans = str.maketrans(text.lower(), repl.lower())
    return text.translate(trans)
assert f(??) == 'lwwer case'

# GPT4 CoT: 'lower case', 'ow'
# could be 'lower case', 'lwwer case'
\end{lstlisting}
\end{minipage}\hfill
\begin{minipage}[t]{.47\textwidth}
\vspace{0pt}
\begin{lstlisting}[caption=GPT-4 CoT input,label={lst:gpt-4-input-fail2}, breaklines=true, language=python]
def f(text):
    string = ''
    for char in text:
        string += char + char.lower()
    return string
assert f(??) == 'llaallaakk'
# GPT-4 CoT: 'LAK'
# should be 'lalak'
\end{lstlisting}
\end{minipage}
\endgroup


\textbf{Other GPT-4 Failures}: Finally, we conclude with a set of six relatively simple string manipulation tasks that we discovered that GPT-4 fails on. We suspect the errors could be partially due to tokenization.
These examples are enough for the thesis point: even the strongest model we tested can fail consistently on small execution problems whose semantics are unambiguous.

\begin{lstlisting}
- What is a string containing 'a' three times, 'b' three times, 'c' twice, 'd' three times, and 'z' twice?
- In Python, what is " BaB ".rfind(" B ")?
- In Python, if I have a string s = 'iabnm~~~~~~~~~~', what is s[1::2]?
- In Python, what is "+".join(['*', '+', 'n', 'z', 'o', 'h'])?
- In Python, if text = "!123Leap and the net will appear" and res = 123, what is text[len(str(res)):]?
- What is "pomodoro".replace("or", "pomodoro")?
\end{lstlisting}



\section{Counterfeit Conundrum: The Nature of Counterfeit Samples}
\label{sec:counterfeit-beyond-score}
\sourcepaper{\path{resources/Counterfeit_Conundrum.zip}. Local source lists Alex Gu as first author, with an equal-contribution author-order note.}{Final thesis prose should be rewritten in dissertation style. If importing close paper phrasing during drafting, wrap it in \texttt{\string\needsrewrite\{...\}} until Alex rewrites it.}
\subsection{Do counterfeit samples from different models or problems differ significantly?} \label{sec:analyzing-counterfeit}
\subsubsection{Is it easier for models to understand counterfeit samples from problems it finds easier?}
Intuitively, if a given programming problem is easy for a model to solve (meaning it has a high pass@1 rate), we might believe models understand how to solve that problem. If that is the case, then models should be able to better understand both correct and counterfeit samples for that problem. To test if this was true, we bucketed problem difficulties into easy, medium, and hard by a model's pass@1 on that problem. We then calculated the average correctness checking, verification, and repair accuracy for each of the tasks for each problem difficulty bucket. Fig. \ref{fig:difficulty-sample} shows a representative subset of these results on HumanEval.

\begin{figure}[ht!]
    \centering
    \begin{subfigure}[b]{0.31\textwidth}
        \centering
        \includegraphics[width=\textwidth]{figures_v2/main/difficulty_barplots/humaneval_correctness_sample.pdf}
        \subcaption{Correctness Checking}
    \end{subfigure}
    \hfill
    \begin{subfigure}[b]{0.31\textwidth}
        \centering
        \includegraphics[width=\textwidth]{figures_v2/main/difficulty_barplots/humaneval_execution_sample.pdf}
        \subcaption{Execution}
    \end{subfigure}
    \hfill
    \begin{subfigure}[b]{0.31\textwidth}
        \centering
        \includegraphics[width=\textwidth]{figures_v2/main/repair/humaneval/repair_4x4.pdf}
        \subcaption{Repair}
    \end{subfigure}
    \caption{Mean accuracy of correctness checking, execution, and repair on HumanEval (error bars are 1 std). A problem's difficulty is determined by the pass@1 of the evaluation model. In general, correctness checking accuracy is uncorrelated with problem difficulty, while both execution and repair accuracies are only very weakly correlated.}
    \label{fig:difficulty-sample}
\end{figure}

Surprisingly, we found that 1) correctness checking accuracies were relatively uncorrelated with problem difficulty, while 2) execution ability and the success rate of repair exhibited a modest amount of correlation with problem difficulty.
We found that these trends were generally robust across HumanEval, LeetCode, and ODEX, although the amount of correlation exhibited in the repair task varied (but was, at best, modest).

To get a more precise measurement of these relationships, we calculated the Pearson correlation between generation performance and each of correctness checking, execution prediction, and repair performance across problems. The resulting histogram showing the distribution of correlations for each task is shown in Fig. \ref{fig:correlation-with-generation}.

\begin{figure}[H]
\centering
\includegraphics[width=0.75\columnwidth]{figures_v2/main/correlation/correlation_with_generation.pdf}
\caption{Correlation between generation and each of our three tasks}
\label{fig:correlation-with-generation}
\end{figure}

\subsubsection{Do models perceive their own samples differently?}
For a given model, its counterfeit samples had a high enough log-likelihood to be generated by the model, so one may hypothesize that models might have a harder time than other models at distinguishing their own counterfeit samples. In Fig. \ref{fig:self-understanding-heatmaps}, we plotted heatmaps showing the performance of various models on datasets generated by other models for the correctness checking task (left) and execution prediction task for HumanEval (right). For both tasks, the relative performance of different models was similar across datasets, and we found no evidence that models faltered more on their own samples. This suggests that counterfeit samples may be general: those from one model are generally difficult for other models to understand as well.
% It would have been really nice to do the perplexity experiment :(

\begin{figure}[H]
    \centering
    \hspace*{\fill}
    \begin{subfigure}[b]{0.4\columnwidth}
        \centering
        \includegraphics[width=\textwidth]{figures_v2/main/accuracy_heatmaps/odex_verification.pdf}
        \subcaption{Correctness (ODEX)}
    \end{subfigure}
    \hspace*{\fill}
    \begin{subfigure}[b]{0.4\columnwidth}
        \centering
        \includegraphics[width=\textwidth]{figures_v2/main/accuracy_heatmaps/humaneval_execution.pdf}
        \subcaption{Execution (HumanEval)}
    \end{subfigure}
    \hspace*{\fill}
    \caption{Heatmap of accuracies for correctness checking and execution prediction. Models do not seem to perceive their own generations differently from those of other models.}
    \label{fig:self-understanding-heatmaps}
\end{figure}

\subsubsection{Do stronger models generate harder counterfeit samples?}
One might also expect that counterfeit samples of stronger models are harder to verify than those of weaker models, as stronger models are less likely to generate obvious mistakes. In Fig. \ref{fig:average-counterfeit-accuracy}, we compared the average scores of two tasks on counterfeit samples for datasets generated by stronger (DS-33B, CL-34B) and weaker (CL-7B, SC) models. Note that this is the same as Fig. \ref{fig:self-understanding-heatmaps} with each row aggregated and filtered to only include counterfeit samples. Since there did not seem to be a significant difference between the difficulties, models of all strengths could be used to generate counterfeit samples that were challenging for models to understand.

\begin{figure}[ht!]
    \centering
    \begin{subfigure}[b]{0.49\columnwidth}
        \centering
        \includegraphics[width=\textwidth]{figures_v2/main/average_model_accuracy/odex_verification.pdf}
        \subcaption{Correctness (ODEX)}
    \end{subfigure}
    \hfill
    \begin{subfigure}[b]{0.49\columnwidth}
        \centering
        \includegraphics[width=\textwidth]{figures_v2/main/average_model_accuracy/humaneval_execution.pdf}
        \subcaption{Execution (HumanEval)}
    \end{subfigure}
    \caption{Accuracy of correctness and execution tasks on counterfeit samples. Both strong and weak models are able to generate difficult counterfeit samples.}
    \label{fig:average-counterfeit-accuracy}
\end{figure}

\subsection{Qualitative Analysis} \label{sec:qual-analysis}
In general, most counterfeit programs fall into one of three broad categories: (1) error in algorithmic design or implementation, (2) incorrectly understanding or completely ignoring details in the specification, and (3) failing to address corner cases in the input space.

\subsubsection{Failure Modes on Correctness Checking}
In Sec. \ref{sec:understanding-counterfeit}, we saw that GPT-4 was significantly more performant than other models at both correctness checking and execution prediction. Yet, we still found a significant fraction of failures. In this section, we took a closer look at the performance of GPT-3.5 and GPT-4 to understand where these models still fell short. Through a manual inspection of examples, we uncovered three main failure modes for verifying the correctness of counterfeit samples and present one example of each. For conciseness and readability, some of the problem statements have been altered from their original form.

%40 manually inspected random samples, 20 per GPT model%, evenly split between the LCB and HumanEval dataset (if yall are concerned that i didnt look at Odex, lmk. I didn't think it was a big deal, and those are much harder to validate manually tbh
% , are classified according to these categories and summarized in Table~\ref{tab:qual_classify}.
% Specific annotated examples are shown in Appendix~\ref{appendix:qual_samples}.

\textbf{1) The model does not catch misunderstood or ignored details in the specification}: Sometimes, the verification model doesn't catch important specification details that are misunderstood or ignored by the counterfeit sample.

% For example, GPT-4 fails to realize that the code below does not actually check that all elements in the subarray are equal, which is explicitly requested in the specification.

% \begin{lstlisting}[language=python]
% """You are given a 0-indexed integer array nums and an integer k. A subarray is called equal if all of its elements are equal. Note that the empty subarray is an equal subarray. Return the length of the longest possible equal subarray after deleting at most k elements from nums. A subarray is a contiguous, possibly empty sequence of elements within an array.
% """

% def longestEqualSubarray(nums: List[int], k: int) -> int:
%     count = defaultdict(int)
%     max_freq = max_length = left = 0
%     for right in range(len(nums)):
%         count[nums[right]] += 1
%         max_freq = max(max_freq, count[nums[right]])
%         if right - left + 1 - max_freq > k:
%             count[nums[left]] -= 1
%             left += 1
%         max_length = max(max_length, right - left + 1)
%     return max_length
% \end{lstlisting}

For example, GPT-3.5 failed to find the reasoning error in the below counterfeit sample because it appeared to have misunderstood the definition of ``corresponding closing bracket''. However, when asked to implement a function for the specification, it correctly understood that the corresponding closing bracket must come after the opening bracket, highlighting an inconsistency between generation and correctness checking.

\begin{lstlisting}[language=python, frame=single, ]
def correct_bracketing(brackets: str):
    """ brackets is a string of "<" and ">". return True if every opening bracket has a corresponding closing bracket.
    >>> correct_bracketing("<<><>>")
    True
    >>> correct_bracketing("><<>")
    False
    """
    opening = brackets.count("<")
    closing = brackets.count(">")
    return opening == closing
\end{lstlisting}

\begin{lstlisting}[style=reasoning, frame=single]
GPT-3.5: The code appears to be correct. It counts the number of opening brackets ("<") and closing brackets (">") in the input string and returns True if the counts are equal, indicating that every opening bracket has a corresponding closing bracket.
\end{lstlisting}

\textbf{2) The model fails to catch subtle implementation mistakes}: Often, counterfeit samples follow a generally correct algorithm design but contain a subtle implementation mistake that models do not catch. Below, the \verb|if num = target| check should consider whether the remaining digits in \verb|square[index:]| sum to 0 for this implementation to be correct, but GPT-4 missed this subtle detail.

\begin{lstlisting}[language=python, frame=single]
"""Given a positive integer n, return the punishment number of n, defined as the sum of the squares of all integers i such that 1) 1 <= i <= n and 2) The decimal representation of i * i can be partitioned into contiguous substrings such that the sum of the integer values of these substrings equals i."""

def punishmentNumber(n: int) -> int:
    def check(num, target, index, square):
        if num == target:
            return True
        if index == len(square):
            return False
        cur = 0
        for i in range(index, len(square)):
            cur = cur * 10 + int(square[i])
            if cur > target:
                break
            if check(num + cur, target, i + 1, square):
                return True
        return False

    res = 0
    for i in range(1, n + 1):
        square = str(i * i)
        if check(0, i, 0, square):
            res += i * i
    return res

# GPT-4 CoT verdict: Correct
\end{lstlisting}

\textbf{3) The model forgives an identified error}: A third failure mode occurs when the model correctly reasons that the code is incorrect, but then makes up a justification for the error and judges the code as correct. In the counterfeit sample below, the code fails to consider capital consonants. While GPT-4 correctly identifies this issue, it proceeds to state that assuming that concern is not a real issue, the code is correct. However, the implementation fails one of the examples.

\begin{lstlisting}[language=python, frame=single]
def select_words(s, n):
    """Given a string s and a natural number n, return a list of all words from string s that contain exactly n consonants, in the order these words appear in the string s.
    Examples:
    select_words("Mary had a little lamb", 4) ==> ["little"]
    select_words("Mary had a little lamb", 3) ==> ["Mary", "lamb"]
    """
    # convert string to list
    word_list = s.split()
    consonants = "bcdfghjklmnpqrstvwxyz"
    # filter out all words that contain n consonants
    return [word for word in word_list if sum(c in consonants for c in word) == n]
\end{lstlisting}

\begin{lstlisting}[style=reasoning, frame=single]
GPT-4: ... However, there is a potential issue with case sensitivity. The `consonants` string only contains lowercase consonants, and the program does not explicitly convert the input string `s` or the words to lowercase before checking if a character is a consonant. This means that if the input string contains uppercase letters, the program might not count them as consonants, leading to incorrect results. [...] the!* program correctly implements the functionality described in the docstring under the assumption that the input string contains only lowercase letters or that case sensitivity is not a concern*!...
\end{lstlisting}


% \begin{lstlisting}[language=python]
% def pairs_sum_to_zero(l):
%     """ pairs_sum_to_zero takes a list of integers as an input. it returns True if there are two distinct elements in the list that sum to zero, and False otherwise.
%     >>> pairs_sum_to_zero([2, 4, -5, 3, 5, 7])
%     True
%     """
%     result = False
%     for i in l:
%         for j in l:
%             if i != j and (i + j == 0):
%                 return True
%     return result
% \end{lstlisting}

% \begin{lstlisting}[style=reasoning]
% Reasoning: ... However, the condition `i != j` checks if the values of `i` and `j` are not equal, not if they are distinct elements (i.e., not checking if they are at different positions in the list). This could lead to a false negative in cases where there are two elements with the same value that could sum to zero with other elements. [...] the program correctly implements the logic as per the docstring for the provided examples and the general case where two distinct elements (by value and position) sum to zero. The misunderstanding comes from interpreting "distinct elements" as distinct by value only, not by position. Since the docstring does not explicitly clarify this point, and based on the examples provided, the program behaves as expected for the specified task.
% \end{lstlisting}


% In order of bottleneck error, t
% The major error classes are as follows. \textbf{Underspecification} happens if the natural language specification does not explicitly state how to handle input corner cases or otherwise does not explicitly state some detail that is important for correctness. \textbf{Misunderstanding the specification}, including fully ignoring some explicitly stated detail, can also result in erroneous verification. In other counterfeit samples, the implementation approach seems to understand the specification just fine but makes an error in either the high-level \textbf{algorithm design} or an \textbf{implementation detail}, and the correctness check verification model does not identify the error.
% In other cases, the \textbf{reasoning and final classification do not align}-- that is, model correctly reasons about the error in an implementation but then ultimately classifies the counterfeit sample as correct.
% Another case is samples that are \textbf{false counterfeit samples}, that is, they are classified as correct by the gold correctness checks, i.e. pass all hidden input-output cases, but they are actually algorithmically incorrect, and the verification model correctly identifies this. % <-- another case of not really a counterfeit sample, but auto methods would have classified it as such so i think we should include it
% also observe that GPT4 reasoning is always longer lol


% \begin{table}[t]
% \small \centering
% \begin{tabular}{lccccc}
% \toprule
%  & Underspecified & Misunderstood spec. & Algo. error & Impl. error  & Concludes adversely \\
% \midrule
% GPT 3.5 & 5 & 8 & 2 & 5 & 0 \\
% GPT 4 & 8 & 3 & 2 & 4 & 3  \\
% \bottomrule
% \end{tabular}
% \caption{40 total counterfeit samples classified according to code correctness verification error. A given counterfeit sample is classified according to its one bottleneck error (left to right).}
% \label{tab:qual_classify}
% \end{table}


% leave everything below this line to future work
% Most errors in correctness classification for counterfeit samples are in misunderstanding the specification and implementation errors. One method to help with underspecification and misunderstanding the specification is to include a correct implementation in the prompt, to provide additional correctness specification signal to the model.
% Overlooking implementation details may be remedied by including the provided counterfeit sample's input/output execution examples in the prompt. If the verification model sees that the output of the code does not align with what would be correct according to the specification, then this may help the model correctly classify the counterfeit sample as ``Incorrect." % We run additional experiments on these hypotheses and show the results in __. We find that indeed both approaches do help improve verification scores, though there is much future work to make that improvement substantial.



% \paragraph{Do the models perform better if given input-output execution of the code?} -- YES slightly



% \paragraph{Do the models perform better if given a ground truth solution of the code?} -- YES

\textbf{Execution Prediction}: To conclude this section, we highlighted an error in execution prediction. In the example below, even though the statement \texttt{if ")" in brackets} is clearly true, GPT-4 was biased by the intended functionality of the program and did not follow its actual execution semantics.

\begin{lstlisting}[language=Python, frame=single]
def correct_bracketing(brackets: str):
    if ")" in brackets:
        return False
    open_brackets = 0
    for bracket in brackets:
        if bracket == "(":
            open_brackets += 1
        else:
            open_brackets -= 1
            if open_brackets < 0:
                return False
    return open_brackets == 0
assert correct_bracketing('()') == ??
# Correct Answer: False
\end{lstlisting}
\begin{lstlisting}[style=reasoning, frame=single]
GPT-4 Execution: [...] 2. The function is called with the argument "()".
!* 3. The first if condition checks if ")" is in brackets. Since it is, but only after "(", this condition does not lead to a return of False at this point [...]
\end{lstlisting}


\subsection{Discussion}
Once again, we saw two surprising nuances that were not clear from benchmark scores. First, we saw a
low correlation between problem difficulty and counterfeit understanding ability. Even for problems that
models pass often, we still found many samples they failed to understand. Second, we discovered that weaker models could generate counterfeits that were as difficult to understand
as stronger models. Overall, counterfeits were a relatively general phenomenon, not just a failure mode
of weaker models or harder problems.


\section{LiveCodeBench: Overfitting, Post-Training, and Open vs. Closed Models}
\label{sec:livecodebench-beyond-score}
% notes
\providecommand{\redo}[1]{{\color{red} #1}}
\providecommand{\rebuttal}[1]{{#1}}


\providecommand{\nbc}[3]{
 {\colorbox{#3}{\bfseries\sffamily\scriptsize\textcolor{white}{#1}}}
 {\textcolor{#3}{\sf\small\textit{#2}}}
 }
\definecolor{annotcolo}{rgb}{0.1,0.1,0.9}
\providecommand{\nj}[1]{\nbc{NJ}{#1}{annotcolo}}


% config
\providecommand{\smalltextsc}[1]{\textsc{\small #1}}
\providecommand{\scripttextsc}[1]{\textsc{\scriptsize #1}}
\providecommand{\tinytextsc}[1]{\textsc{\tiny #1}}

% sysnames
\providecommand{\livecodebench}{\smalltextsc{LiveCodeBench}}
\providecommand{\lcb}{\smalltextsc{LCB}}
\providecommand{\livecodebenchtitle}{\textsc{LiveCodeBench}}


% general
\providecommand{\llm}{\smalltextsc{LLM}\xspace}
\providecommand{\llms}{\smalltextsc{LLMs}\xspace}
\providecommand{\chainot}{\smalltextsc{COT}\xspace}
\providecommand{\python}{\smalltextsc{Python}\xspace}




%models

%openai
\providecommand{\openai}{\smalltextsc{OpenAI}}
\providecommand{\gpts}{\smalltextsc{GPTs}\xspace}
\providecommand{\codedavinci}{\smalltextsc{Code-Davinci-002}\xspace}
\providecommand{\gptthreefiveturbo}{\smalltextsc{GPT-3.5-turbo}\xspace}
\providecommand{\gptfour}{\smalltextsc{GPT-4}\xspace}
\providecommand{\gptfourturbo}{\smalltextsc{GPT-4-Turbo}\xspace}
\providecommand{\gptfouro}{\smalltextsc{GPT-4-O}\xspace}

%deepmind
\providecommand{\alphacode}{\smalltextsc{AlphaCode}\xspace}
\providecommand{\alphacodeB}[1]{\smalltextsc{AlphaCode-#1B}\xspace}

%anthropic
\providecommand{\claude}{\smalltextsc{Claude}\xspace}
\providecommand{\claudes}{\smalltextsc{Claudes}\xspace}
\providecommand{\claudetwo}{\smalltextsc{Claude-2}\xspace}
\providecommand{\claudethrees}{\smalltextsc{Claude-3s}\xspace}
\providecommand{\claudeinstantone}{\smalltextsc{Claude-Ins-1}\xspace}
\providecommand{\claudeopus}{\smalltextsc{Claude-3-Opus}\xspace}
\providecommand{\claudesonnet}{\smalltextsc{Claude-3-Sonnet}\xspace}

%google
\providecommand{\gemini}{\smalltextsc{Gemini}\xspace}
\providecommand{\geminis}{\smalltextsc{Geminis}\xspace}
\providecommand{\geminipro}{\smalltextsc{Gemini-Pro}\xspace}
\providecommand{\geminiproonefive}{\smalltextsc{Gemini-Pro-1.5}\xspace}
\providecommand{\geminiflash}{\smalltextsc{Gemini-Flash}\xspace}
\providecommand{\geminiultra}{\smalltextsc{Gemini-Ultra}\xspace}

%meta 
\providecommand{\cllamas}{\smalltextsc{CodeLLaMas}\xspace}
\providecommand{\cllama}{\smalltextsc{CodeLLaMa}\xspace}
\providecommand{\cllamains}{\smalltextsc{CodeLLaMa-Ins}\xspace}
\providecommand{\cllamabase}{\smalltextsc{CodeLLaMa-Base}\xspace}
\providecommand{\cllamaB}[1]{\smalltextsc{CL-Ins-#1B}\xspace}
\providecommand{\cllamacommonB}[1]{CL-#1B}
\providecommand{\cllamabaseB}[1]{\smalltextsc{CL-Base-#1B}\xspace}
\providecommand{\cllamaBList}{\smalltextsc{CL-Ins-\{7, 13, 34\}B}\xspace}
\providecommand{\cllamabaseBList}{\smalltextsc{CL-Base-\{7, 13, 34\}B}\xspace}

\providecommand{\llamas}{\smalltextsc{LLaMa-3s}\xspace}

\providecommand{\llamabase}{\smalltextsc{L3-Base}}
\providecommand{\llamabaselist}{\smalltextsc{L3-Base-\{7, 70\}B}\xspace}
\providecommand{\llamainslist}{\smalltextsc{L3-Ins-\{7, 70\}B}\xspace}
\providecommand{\llamaB}[1]{\smalltextsc{L3-Ins-#1B}\xspace}
\providecommand{\llamabaseB}[1]{\smalltextsc{L3-Base-#1B}\xspace}

%finetunes
\providecommand{\wizardlong}{\smalltextsc{WizardCoder}\xspace}
\providecommand{\wizard}{\smalltextsc{WC-34B}\xspace}
\providecommand{\wizardtwo}{\smalltextsc{WC-33B}\xspace}
\providecommand{\phind}{\smalltextsc{Phind-34B}\xspace}
\providecommand{\magicoders}{\smalltextsc{MagiCoders}\xspace}
\providecommand{\magicoder}[1]{\smalltextsc{MC-#1B}\xspace}
\providecommand{\magicoderBLit}{\smalltextsc{MC-\{6.7, 7\}B}\xspace}

%deepseek
\providecommand{\deepseek}{\smalltextsc{DeepSeek}\xspace}
\providecommand{\deepseekbase}{\smalltextsc{DeepSeek-Base}\xspace}
\providecommand{\deepseekinstruct}{\smalltextsc{DeepSeek-Instruct}\xspace}
\providecommand{\deepseeks}{\smalltextsc{DeepSeeks}\xspace}
\providecommand{\deepseekcode}[1]{\smalltextsc{DS}-#1B\xspace}
\providecommand{\deepseekcodeins}{\smalltextsc{DS-Ins}\xspace}
\providecommand{\deepseekcodeB}[1]{\smalltextsc{DS-Ins-#1B}\xspace}
\providecommand{\deepseekbaseB}[1]{\smalltextsc{DS-Base-#1B}\xspace}
\providecommand{\deepseekcodeBList}{\smalltextsc{DS-Ins-\{1.3, 6.7, 33\}B}\xspace}
\providecommand{\deepseekcodebaseBList}{\smalltextsc{DS-Base-\{1.3, 6.7, 33\}B}\xspace}

% codeqwen
\providecommand{\codeqwen}{\smalltextsc{CodeQwen}\xspace}
\providecommand{\codeqwenchat}{\smalltextsc{CodeQwen-Chat}\xspace}

%mistral

\providecommand{\mixtral}{\smalltextsc{Mixtral}\xspace}
\providecommand{\codestral}{\smalltextsc{Codestral}\xspace}
\providecommand{\mistrallarge}{\smalltextsc{Mistral-L}\xspace}
\providecommand{\mistral}{\smalltextsc{Mistral}\xspace}

% bigcode

\providecommand{\starcoder}{\smalltextsc{StarCoder}\xspace}
\providecommand{\starcodertwo}{\smalltextsc{StarCoder2}\xspace}
\providecommand{\starcodertwobase}{\smalltextsc{StarCoder2-Base}\xspace}
\providecommand{\sctwoinsB}[1]{\smalltextsc{SC2-#1B}\xspace}
\providecommand{\sctwobaseB}[1]{\smalltextsc{SC2-Base-#1B}\xspace}
\providecommand{\sctwobaseBList}{\smalltextsc{SC2-Base-\{3,7,15\}B}\xspace}

% cohere
\providecommand{\commandr}{\smalltextsc{CMD-R}\xspace}
\providecommand{\commandrplus}{\smalltextsc{CMD-R+}\xspace}


% datasets
\providecommand{\thestack}{\smalltextsc{Stack}\xspace}
\providecommand{\apps}{\smalltextsc{APPS}\xspace}
\providecommand{\contests}{\smalltextsc{Code-Contests}\xspace}
\providecommand{\humaneval}{\smalltextsc{HumanEval}}
\providecommand{\humanevalplus}{\smalltextsc{HumanEval+}}
\providecommand{\mbpp}{\smalltextsc{MBPP}}
\providecommand{\cruxeval}{\smalltextsc{CRUXEval}}
\providecommand{\codet}{\smalltextsc{CodeT}}
\providecommand{\xcode}{\smalltextsc{xCodeEval}\xspace}
\providecommand{\codescope}{\smalltextsc{CodeScope}\xspace}
\providecommand{\taco}{\smalltextsc{TACO}\xspace}
\providecommand{\usacobench}{\smalltextsc{USCAOBench}\xspace}
%coding websites
\providecommand{\leetcode}{\smalltextsc{LeetCode}\xspace}
\providecommand{\atcoder}{\smalltextsc{AtCoder}\xspace}
\providecommand{\codeforces}{\smalltextsc{CodeForces}\xspace}
\providecommand{\html}{\smalltextsc{HTML}\xspace}

\providecommand{\easy}{\smalltextsc{Easy}\xspace}
\providecommand{\medium}{\smalltextsc{Medium}\xspace}
\providecommand{\hard}{\smalltextsc{Hard}\xspace}

%website scrapes
\providecommand{\problem}{\smalltextsc{$P$}\xspace}
\providecommand{\solution}{\smalltextsc{$S$}\xspace}
\providecommand{\tests}{\smalltextsc{$T$}\xspace}
\providecommand{\testspublic}{\smalltextsc{$T_{p}$}\xspace}
\providecommand{\testshidden}{\smalltextsc{$T_{h}$}\xspace}
\providecommand{\contestdate}{\smalltextsc{$D$}\xspace}

% metrics
\providecommand{\passmetric}[1]{\smalltextsc{Pass@}#1\xspace}



% %%%%
\lstset{
  basicstyle=\ttfamily\small,
  columns=fullflexible,
  % frame=single,
  breaklines=true,
  postbreak=\mbox{$\hookrightarrow$\space},
  rulecolor=\color{black},
}
\definecolor{codegreen}{rgb}{0,0.6,0}
\lstdefinestyle{mystyle}{  
    commentstyle=\color{codegreen},
    keywordstyle=\color{blue},
    basicstyle=\ttfamily\small,
    breakatwhitespace=false,        
    breaklines=true,                 
    captionpos=b,                    
    keepspaces=true,                 
    showspaces=false,                
    showstringspaces=false,
    showtabs=false,                  
    tabsize=2
}
\lstset{style=mystyle}

% if you use cleveref..

% %%%%%%%%%%%%%%%%%%%%%%%%%%%%%%%%
% % THEOREMS
% %%%%%%%%%%%%%%%%%%%%%%%%%%%%%%%%
% \theoremstyle{plain}
% \newtheorem{theorem}{Theorem}[section]
% \newtheorem{proposition}[theorem]{Proposition}
% \newtheorem{lemma}[theorem]{Lemma}
% \newtheorem{corollary}[theorem]{Corollary}
% \theoremstyle{definition}
% \newtheorem{definition}[theorem]{Definition}
% \newtheorem{assumption}[theorem]{Assumption}
% \theoremstyle{remark}
% \newtheorem{remark}[theorem]{Remark}

% % Todonotes is useful during development; simply uncomment the next line
% %    and comment out the line below the next line to turn off comments

% \input{math_commands.tex}
% \newcommand{\todo}[1]{\textcolor{red}{[TODO: #1]}}
\newcommand{\term}[1]{\emph{#1}}
\definecolor{rewriteblue}{RGB}{0,70,170}
\definecolor{mygreen}{RGB}{226,240,217}
\definecolor{myblue}{RGB}{222,235,247}
\definecolor{myorange}{RGB}{251,229,214}
\definecolor{lightblue}{RGB}{232,242,255}
\definecolor{lightorange}{RGB}{255,241,224}
\definecolor{good}{HTML}{1B7837}
\definecolor{bad}{HTML}{B2182B}
\newcommand{\needsrewrite}[1]{\textcolor{rewriteblue}{[SOURCE-DERIVED; REWRITE: #1]}}
\newenvironment{rewriteblock}{%
  \par\begingroup\color{rewriteblue}%
  \noindent
}{%
  \par\endgroup
}
\newenvironment{draftblock}{%
  \par\noindent
}{%
  \par
}
\newcommand{\sourcepaper}[2]{}

\definecolor{lightgreen}{HTML}{C4FFCF}
\definecolor{codegreen}{rgb}{0,0.6,0}

\newcommand{\red}[1]{\textbf{\textcolor{red}{#1}}}
\newcommand{\green}[1]{\textbf{\textcolor{ForestGreen}{#1}}}
\newcommand{\cmark}{\ding{51}}
\newcommand{\xmark}{\ding{55}}
\newcommand{\benchmark}{\textsc{CRUXEval}\xspace}
\newcommand{\benchmarki}{\textsc{CRUXEval-I}\xspace}
\newcommand{\benchmarko}{\textsc{CRUXEval-O}\xspace}
\newcommand{\codellamamed}{\textsc{Code Llama 13B}\xspace}
\newcommand{\codellamalarge}{\textsc{Code Llama 34B}\xspace}
\providecommand{\multirow}[3]{#3}
\providecommand{\hdashline}[1][]{\midrule}
\NewDocumentCommand{\cdashline}{m o}{}
\providecommand{\makecell}[1]{\begin{tabular}{@{}c@{}}#1\end{tabular}}
\providecommand{\smallsec}[1]{\textbf{#1\,\,}}
\providecommand{\llm}{\textsc{LLM}\xspace}
\providecommand{\llms}{\textsc{LLMs}\xspace}
\providecommand{\lean}{\textsc{Lean}\xspace}
\providecommand{\setstretch}[1]{}
\providecommand{\smallsim}{\mathrel{\sim}}
\providecommand{\subfloat}[2][]{%
  \begin{minipage}{0.49\textwidth}%
  \centering
  #2%
  \if\relax\detokenize{#1}\relax\else\subcaption{#1}\fi
  \end{minipage}%
}

\DeclareMathOperator*{\argmax}{arg\,max}
\DeclareMathOperator*{\argmin}{arg\,min}
\lstdefinelanguage{Lean}{
  morekeywords={theorem,lemma,def,by,have,show,calc,let,fun,match,with,if,then,else,import,open,namespace,end,where,structure,inductive,example},
  sensitive=true,
  morecomment=[l]{--},
  morecomment=[s]{/-}{-/},
  morestring=[b]"
}
\lstdefinelanguage{triton}{
  morekeywords={tl,triton,jit,program_id,load,store,arange,zeros,where,constexpr,block_ptr,dtype},
  sensitive=true,
  morecomment=[l]{\#},
  morestring=[b]"
}
\lstdefinelanguage{CUDA}{
  morekeywords={__global__,__device__,__shared__,__constant__,__host__,__syncthreads,threadIdx,blockIdx,blockDim,gridDim},
  sensitive=true,
  morecomment=[l]{//},
  morecomment=[s]{/*}{*/},
  morestring=[b]",
  morestring=[b]'
}
\lstdefinelanguage{Rust}{
  morekeywords={as,break,const,continue,crate,else,enum,extern,false,fn,for,if,impl,in,let,loop,match,mod,move,mut,pub,ref,return,self,Self,static,struct,super,trait,true,type,unsafe,use,where,while,async,await,dyn},
  sensitive=true,
  morecomment=[l]{//},
  morecomment=[s]{/*}{*/},
  morestring=[b]",
  morestring=[b]'
}
\lstdefinelanguage{Python}{
  morekeywords={and,as,assert,break,class,continue,def,del,elif,else,except,False,finally,for,from,global,if,import,in,is,lambda,None,nonlocal,not,or,pass,raise,return,True,try,while,with,yield},
  sensitive=true,
  morecomment=[l]{\#},
  morecomment=[s]{\"\"\"}{\"\"\"},
  morecomment=[s]{\'\'\'}{\'\'\'},
  morestring=[b]',
  morestring=[b]"
}
\lstdefinelanguage{json}{
  morekeywords={true,false,null},
  sensitive=false,
  morestring=[b]",
  morecomment=[l]{//},
}

\lstdefinestyle{mystyle}{
  commentstyle=\color{codegreen},
  keywordstyle=\color{blue},
  basicstyle=\ttfamily\footnotesize,
  breakatwhitespace=false,
  breaklines=true,
  captionpos=b,
  keepspaces=true,
  showspaces=false,
  showstringspaces=false,
  showtabs=false,
  tabsize=2
}
\lstdefinestyle{boxed}{
  frame=single,
  breaklines=true,
  postbreak=\mbox{$\hookrightarrow$\space},
  commentstyle=\color{codegreen},
  keywordstyle=\color{blue},
  basicstyle=\ttfamily\footnotesize,
  breakatwhitespace=false,
  captionpos=t,
  keepspaces=true,
  showspaces=false,
  showstringspaces=false,
  showtabs=false,
  tabsize=2
}
\lstdefinestyle{specblock}{
  frame=single,
  breaklines=true,
  basicstyle=\ttfamily\scriptsize,
  keepspaces=true,
  showstringspaces=false,
  columns=fullflexible
}
\DefineVerbatimEnvironment{codebox}{Verbatim}{fontsize=\small}
\DefineVerbatimEnvironment{examplecodebox}{Verbatim}{fontsize=\small}
\lstdefinestyle{reasoning}{
  breaklines=true,
  commentstyle=\color{codegreen},
  keywordstyle=\color{blue},
  breakautoindent=0pt,
  breakindent=0pt,
  basicstyle=\itshape\footnotesize,
  breakatwhitespace=true,
  captionpos=t,
  keepspaces=true,
  showspaces=false,
  showstringspaces=false,
  showtabs=false
}
\lstdefinestyle{mystylesmaller}{
  commentstyle=\color{codegreen},
  keywordstyle=\color{blue},
  basicstyle=\ttfamily\scriptsize,
  breaklines=true,
  captionpos=b,
  keepspaces=true,
  showspaces=false,
  showstringspaces=false,
  showtabs=false,
  tabsize=2
}
\lstdefinestyle{mystyletiny}{
  commentstyle=\color{codegreen},
  keywordstyle=\color{blue},
  basicstyle=\ttfamily\tiny,
  breaklines=true,
  captionpos=b,
  keepspaces=true,
  showspaces=false,
  showstringspaces=false,
  showtabs=false,
  tabsize=2
}
\lstdefinestyle{leancompact}{
  language=Lean,
  basicstyle=\fontsize{3}{3.5}\ttfamily,
  breaklines=true,
  captionpos=b,
  keepspaces=true,
  showspaces=false,
  showstringspaces=false,
  showtabs=false,
  frame=none
}
\lstset{
  style=mystyle,
  columns=fullflexible,
  frame=single,
  postbreak=\mbox{$\hookrightarrow$\space},
  rulecolor=\color{black},
  moredelim=[is][\color{red}]{!*}{*!},
  literate={₀}{{$_0$}}1
           {₁}{{$_1$}}1
           {₂}{{$_2$}}1
           {₃}{{$_3$}}1
           {₄}{{$_4$}}1
           {₅}{{$_5$}}1
           {₆}{{$_6$}}1
           {₇}{{$_7$}}1
           {₈}{{$_8$}}1
           {₉}{{$_9$}}1
           {¹}{{$^1$}}1
           {²}{{$^2$}}1
           {³}{{$^3$}}1
           {⁻}{{$^-$}}1
           {·}{{$\cdot$}}1
           {×}{{$\times$}}1
           {∀}{{$\forall$}}1
           {∃}{{$\exists$}}1
           {ℕ}{{$\mathbb{N}$}}1
           {ℚ}{{$\mathbb{Q}$}}1
           {ℝ}{{$\mathbb{R}$}}1
           {ℤ}{{$\mathbb{Z}$}}1
           {∧}{{$\wedge$}}1
           {∨}{{$\vee$}}1
           {∈}{{$\in$}}1
           {∉}{{$\notin$}}1
           {∑}{{$\sum$}}1
           {∫}{{$\int$}}1
           {∣}{{$\mid$}}1
           {≠}{{$\ne$}}1
           {≡}{{$\equiv$}}1
           {≤}{{$\le$}}1
           {≥}{{$\ge$}}1
           {←}{{$\leftarrow$}}1
           {→}{{$\to$}}1
           {↑}{{$\uparrow$}}1
           {⊢}{{$\vdash$}}1
           {⟨}{{$\langle$}}1
           {⟩}{{$\rangle$}}1
           {…}{{\ldots}}1
}

\makeatletter
\def\@captype{lstlisting}
\lst@AddToHook{Init}{\def\@captype{lstlisting}}
\makeatother

\theoremstyle{plain}
\newtheorem{theorem}{Theorem}[chapter]
\newtheorem{lemma}[theorem]{Lemma}
\newtheorem{proposition}[theorem]{Proposition}
\newtheorem{corollary}[theorem]{Corollary}

\theoremstyle{definition}
\newtheorem{definition}[theorem]{Definition}
\newtheorem{example}[theorem]{Example}

\theoremstyle{remark}
\newtheorem{remark}[theorem]{Remark}



\providecommand{\smallsim}{\smallsym{\mathrel}{\sim}}

\makeatletter
\providecommand{\smallsym}[2]{#1{\mathpalette\make@small@sym{#2}}}
\providecommand{\make@small@sym}[2]{%
  \vcenter{\hbox{$\m@th\downgrade@style#1#2$}}%
}
\providecommand{\downgrade@style}[1]{%
  \ifx#1\displaystyle\scriptstyle\else
    \ifx#1\textstyle\scriptstyle\else
      \scriptscriptstyle
  \fi\fi
}
\makeatother


\subsection[Overfitting to HumanEval+]{Overfitting to HumanEval+}

We first compare the code generation performance of models on \humanevalplus{} vs. \livecodebench{},
shown in Figure~\ref{fig:he_vs_lcb}. We observe two distinct clusters of models, shown in red and green. The green cluster shows models that perform 
similarly on both benchmarks, consisting mostly of base and closed-access models. 
The red cluster shows models that perform significantly better on \humanevalplus{} than \livecodebench{}. Careful analysis reveals that
this cluster consists primarily of fine-tuned variants of open-access models, strongly indicative of overfitting
or data contamination.  In particular, we draw attention to \deepseekcodeB{1.3}, which achieves
$60\%$ on \humanevalplus{} but only $26\%$ on \lcb{}-Easy.
These results mirror what we found with \benchmark{} in Section~\ref{sec:cruxeval-beyond-score}: a class of
models with a high HumanEval (\humanevalplus{} in this case) score with no corresponding gains
on a different benchmark.

Overall, the lesson is that high scores on familiar benchmarks can be ambiguous.
Sometimes, it may indicate general programming abilities, but other times it may be due to benchmark-specific adaptation. 
It is crucial to compare across a varied suite of different benchmarks to understand the full picture, and we should be 
careful of techniques that promote overfitting rather than developing general abilities.

\begin{figure}[!t]
    \centering
    \includegraphics[width=0.99\textwidth]{figures/main_results/lcb_vs_he.png}
  \vspace{-3pt}
    \caption{
        Scatter plot comparing \passmetric{1} of models on \humanevalplus{} versus \lcb{}-Easy{} (time-window Sep $2023$ to May $2024$).
        The green-shaded region marks models with aligned performance across benchmarks, while the red-shaded region marks models with much stronger performance on \humanevalplus{}.
    }

        \label{fig:he_vs_lcb}
        \vspace{-15pt}
\end{figure}


\subsection{Role of Post-Training}
Next, we analyze the performance of base models vs. post-trained models. In terms of 
base models, we see that \llamabase{} and \deepseek{} significantly outperform \cllama{} and
\starcodertwo. However, these differences may partially be due to the training data
distribution. For example, \starcodertwo uses competition coding problems in its
pre-training corpus.

Our benchmark also shows that post-training significantly improves performance, as 
\llamaB{70}, \deepseekcodeB{33}, and \phind{} achieve 28.3, 23.6, and 21.0 \passmetric{1} on \lcb{},
improving over their base models by 8.2, 7.3, and 9.5 points respectively.

On the other hand, while base models generally fall into the green cluster of models,
post-trained models generally end up in the red cluster.  
We hypothesize that compared to data used in closed-access models, the fine-tuning 
datasets for open models are lacking in diversity and quality, which may lead to poorer
generalization. 

\subsection{Gap Between Closed and Open Models}

In Figure~\ref{fig:lcb-beyond-open-closed-gap}, we investigate the gap between
open-access and closed-access models. 
We already started to see this Section~\ref{sec:cruxeval-beyond-score}, where we found that GPT-4 had already 
distinguished itself at CRUXEval compared to other models. It showed good performance on both CRUXEval and HumanEval seemingly
without overfitting to either, and benefited significantly more from chain-of-thought than 
any other models. 
Overall, on \livecodebench{}, closed models generally outperform open models.
However, the gap is not uniform, as models like \llamaB{70}, \mixtral{}, and \deepseekcodeB{33}
perform on-par with the weaker closed models like \geminipro{} and \gptthreefiveturbo{}, and 
\claudesonnet. However, this varies across scenarios, as closed-access models perform better in the 
test output prediction and code execution scenarios. This suggests that strong base models
and high-quality post-training data can partially close the gap, but that there are several
capabilities where closed models still have a significant lead. 


\begin{figure}[H]
  \centering
  \includegraphics[width=0.65\linewidth]{figures/overview/lc_barchart.png}
  \caption{Comparison of open-access and closed-access models on the
  \livecodebench{}-Easy code-generation scenario.}
  \label{fig:lcb-beyond-open-closed-gap}
\end{figure}


\section{LINC: A Neurosymbolic Approach for Logical Reasoning by Combining Language Models with First-Order Logic Provers}
\label{sec:linc-beyond-score}
\sourcepaper{\path{resources/_EMNLP_2023__LINC_Logical_Reasoning.zip}. Local source lists equal contribution with randomized author order; Alex Gu can be treated as first author for thesis framing.}{Because parts of the LINC paper prose were written by coauthors, any source-derived LINC text should still be wrapped in \texttt{\string\needsrewrite\{...\}} until Alex rewrites it in his own words.}

% Verbatim draft import from camera-ready.tex.
% Section levels and labels are adjusted for thesis inclusion under Section 3.5.

In this section, we shifted our focus towards logical reasoning: deductively inferring the truth value of a
conclusion from a collection of premises. Prior to our work, there had been many prompting and chain-of-thought based
strategies to improve the logical reasoning capabilities of language models. However, we found that they
failed in subtle and unpredictable ways. To mitigate this, we introduced \texttt{LINC}: \textbf{L}ogical
\textbf{I}nference via \textbf{N}eurosymbolic \textbf{C}omputation.

In \texttt{LINC}, logical reasoning is tackled through a two-step neurosymbolic process.
In the first step, a language model takes the natural language premises and desired conclusion and converts them into first-order logic (FOL) expressions \cite{enderton2001mathematical, barker-plummer2011language}.
Once everything is represented in first-order logic, we then apply a symbolic theorem prover to determine the truth value of the conclusion.
As the translation from natural language to first-order logic can be imperfect, we use a majority voting scheme to
improve the reliability of the approach.

The key advantage of \texttt{LINC} is that the reasoning burden is offloaded from the language model to the theorem prover.
On the other hand, the challenge shifts to faithful and accurate formalization. The FOL translation must capture all the relevant
information from the premises without adding or removing any information. Since it is not clear how the difficulties of formalization
and natural language reasoning compare, the main focus of our work was to compare and contrast \texttt{LINC} with traditional
\texttt{Chain-of-Thought} approaches. Through thorough experimentation and analysis on two challenging datasets,
we provided insights into the strengths and weaknesses of both approaches. Overall, we found that they have distinct behaviors and
failure modes, suggesting a strong synergy between theorem provers and language models.


\begin{figure}[ht!]
  \centering
  \includegraphics[width=\linewidth]{figures/pipeline.pdf}
  \caption{Overview of the \texttt{LINC} approach. Starting from a problem in natural language, the language model first samples first-order logic translations of the premises and conclusion. The translated formulas are then passed to an automated theorem prover, which filters invalid syntax and predicts labels for the valid samples. Finally, a majority-vote step aggregates the remaining candidate labels.}
  \label{linc-paper:fig:main}
\end{figure}

%%%
\subsection{\texttt{LINC}: Logical Inference via Neurosymbolic Computation}
\label{linc-paper:sec:ns-model}

Fig.~\ref{linc-paper:fig:main} provides an overview of \texttt{LINC}'s two-stage approach.
First, the LLM acts as a semantic parser to generate first-order logic expressions to be sent to the automated theorem prover.
Second, \texttt{Prover9} \citep{prover9-mace4} executes the expressions to determine whether the conclusion is \texttt{True},
\texttt{False}, or \texttt{Uncertain} under the translated premises, and raises an exception if the FOL syntax is
invalid. The $K$-Maj voting scheme then aggregates predictions across multiple samples to determine the final prediction.

The core tradeoff of \texttt{LINC} is that the difficulty is shifted from faithful reasoning to faithful formalization.
Using \texttt{LINC}, a language model no longer needs to reason in natural language. This
removes any reasoning errors that may arise from the model's inability to properly keep
of logical structure. On the other hand, with \texttt{LINC}, the model must now faithfully formalize
each of the natural language statements into FOL. This requires the FOL expressions to both be 1) \textit{syntactically}
valid, so that they are accepted by the prover, and 2) \textit{semantically} valid, so that they contain the same semantic information
as the natural language statements. As shown later in the error analysis, this translation can be difficult when statements are more complex,
because the FOL expressions may fail to capture both explicit and implicit information.

%%%
\subsection{Experiments}\label{linc-paper:sec:experiments}
In this section, we present our experimental setup, the models we used, and the three baselines to which we compared \texttt{LINC}.

\textbf{Datasets:} We used two challenging datasets for logical reasoning in natural language: \textbf{FOLIO} \cite{han2022folio} and \textbf{ProofWriter} \cite{tafjord2021proofwriter}.
FOLIO is an expert-written, open-domain, logically complex and diverse dataset for natural
language reasoning with first-order logic. We used the validation set but discovered that 22 of the 204 samples contain errors
(details in the original paper appendix), so we used the other 182 samples.
ProofWriter is a synthetically generated dataset for logical reasoning over natural language.
We randomly selected 360 samples from the OWA (Open-World Assumption) subset of ProofWriter in
a way that is balanced across the number of reasoning steps in the shortest ground truth proof
(depth 0-5; 50 samples each) and across the three labels (120 samples each; 20 each per depth).

\textbf{In-context learning examples:} We hand-picked eight diverse samples from the FOLIO training set
to be used as few-shot in-context examples. Because ProofWriter does not come with ground truth
FOL statements, we used these eight samples for both evaluations. Questions in
ProofWriter generally have more premises per question than in FOLIO (5.3 in FOLIO vs. 18.8 in ProofWriter).
Thus, our evaluation on FOLIO is an \emph{in-distribution} task, whereas ProofWriter
requires generalizing \emph{out-of-distribution} to reasoning over considerably larger
sets of premises than are given in the prompt.

\textbf{Majority voting:} We used $K$-way majority voting \cite{wang2023self}, where $K$ samples are taken i.i.d. from the model and the mode used as the final prediction.
We reported $K{=}10$-way majority voting unless otherwise stated, with ties broken arbitrarily.
We reported the effect of $K$ on performance across our conditions and briefly discussed trends in the original paper appendix.

\textbf{Models:}
We used three models pre-trained on both natural language and code: GPT-3.5 \citep{ouyang2022training}, GPT-4 (\texttt{gpt-3.5-turbo-16k-0613} and \texttt{gpt-4-0613}) \citep{openai2023gpt4}, and StarCoder+ (\url{https://huggingface.co/bigcode/starcoderplus}) \citep{li2023starcoder} with a decoding temperature of $T=0.8$ for all experiments. We defer model, hyperparameter, and hardware details to the original paper appendix.
We opted for \texttt{StarCoder+} for three reasons: firstly, unlike the other models we considered, it is a free, open-access model. Secondly, it has a dataset search functionality (see \url{https://huggingface.co/spaces/bigcode/in-the-stack} and \url{https://huggingface.co/spaces/bigcode/search}), with which we verified that FOLIO and ProofWriter were not in \texttt{StarCoder+}'s training set, giving further assurance in the validity of our findings. Thirdly, with its 15.5B parameters it is likely considerably smaller than GPT-3.5 and GPT-4; although the exact size of these models has not been made public, their common predecessor GPT-3 was known to have 175B parameters \citep{brown2020language}. This allowed us to compare performance at different model scales.

\textbf{Controlled baselines:}

\begin{figure}[ht!]
  \includegraphics[width=\linewidth]{figures/prompt.pdf}
  \caption{This figure outlines the string concatenation workflow for each of our conditions. We start with the original problem, provide ICL examples through an intermediate markup language, and finally append the problem to evaluate. At this stage, we allow the model to autoregressively sample until producing a stop token.}
  \label{linc-paper:fig:prompt}
\end{figure}
\label{linc-paper:sec:controlled-baselines}
As shown in Fig.~\ref{linc-paper:fig:prompt}, we compared LINC with three baseline approaches: \texttt{Na\"ive}, \texttt{Scratchpad}, and \texttt{Chain-of-Thought (CoT)}.
In \texttt{Na\"ive} (red), the model is given the natural language premises and asked to directly generate the label (\texttt{True}/\texttt{False}/\texttt{Uncertain}).
In \texttt{Scratchpad} (blue) \cite{nye2021show}, we ablated the effectiveness of using a symbolic prover instead of a LLM in the second stage.
The model is asked to generate FOL expressions just like in LINC, but instead of passing the FOL expressions to a symbolic prover, we ask the model to directly predict the label based on its own FOL translation.
In \texttt{CoT} (yellow), we used standard \texttt{CoT} prompting \cite{wei2022chain,kojima2022large,wang2023self}, where the model is asked to generate step-by-step natural language reasoning to arrive at the conclusion.
The prompts we used for all approaches can be found in the original paper appendix.

\begin{figure}[th]
  \centering
  \begin{subfigure}[b]{0.45\textwidth}
    \centering
    \includegraphics[width=\textwidth]{figures/folio_bar_chart.pdf}
    \caption{FOLIO.}
    \label{linc-paper:fig:results-folio}
  \end{subfigure}
  \hfill
  \begin{subfigure}[b]{0.45\textwidth}
    \centering
    \includegraphics[width=\textwidth]{figures/proofwriter_bar_chart.pdf}
    \caption{ProofWriter.}
    \label{linc-paper:fig:results-pw}
  \end{subfigure}
  \caption{Results of each model on the FOLIO and ProofWriter datasets. Accuracies are for bootstrapped $10$-way majority vote for all models. Error bars are $\pm 1$ bootstrapped standard deviation. Dotted, black line is the accuracy obtained by always guessing the most common label in the dataset.}
  \label{linc-paper:fig:results-main}
\end{figure}

%%%
\subsection{Results \& Discussion}
\label{linc-paper:sec:results_discussion}

Our main results are shown in Figure~\ref{linc-paper:fig:results-main}.
Each bar represents either \texttt{LINC} or one of the three baselines,
while each group of bars indicates the language model used (in \{\texttt{StarCoder+}, \texttt{GPT-3.5}, \texttt{GPT-4}\}).

On FOLIO (Figure \ref{linc-paper:fig:results-folio}), StarCoder+---the smallest model we experimented with---benefited the most from \texttt{LINC}, achieving a mean accuracy that was 14.2 points higher than the closest controlled baseline (56.0\% vs. 41.8\% with \texttt{CoT}).
We found that \texttt{Scratchpad}, where the intermediate logical formulae are still generated but the call to the symbolic solver is ablated and replaced by the model's own prediction, did not appear to benefit performance; for StarCoder+, neither \texttt{Scratchpad} nor the \texttt{Na\"ive} baseline performed better than simply deterministically predicting the most common label (``\texttt{Uncertain}'').
For GPT-3.5, the trend is similar, although the gap between \texttt{LINC} and the closest baseline shrinks (62.6\% vs. 54.9\% average accuracies).
For GPT-4, the trend reversed: \texttt{LINC} underperformed \texttt{CoT}. However, we performed a McNemar's test \cite{mcnemar1947note} to get the p-value on this GPT-4 \texttt{LINC} vs. \texttt{CoT} comparison, and we found that the difference was not significant ($p=0.58$).
Meanwhile, for our balanced subset of ProofWriter, we saw significant performance gains across the board (Figure \ref{linc-paper:fig:results-pw}); particularly so for GPT-3.5 and GPT-4, which achieved mean accuracies of 96.4\% and 98.3\% when paired with LINC.

We attribute the superior performance of \texttt{LINC} on ProofWriter to two causes.
First, ProofWriter is much easier to formalize than FOLIO, which is the bottleneck for LINC.
While FOLIO contains complex sentences with linguistic variety, all sentences in ProofWriter are synthetically generated and easily formalizable.
Second, ProofWriter is an \emph{out-of-distribution} setting, where the in-context examples contain fewer premises and shorter deductive chains than the
reasoning task.
This distribution shift makes it harder for the model to ignore irrelevant premises and keep track of the long deductive chains.
By removing the reasoning burden, \texttt{LINC} is able to ignore this difficulty and perform well.

\begin{figure}[th]
  \centering
  \begin{subfigure}[b]{0.32\textwidth}
    \centering
    \includegraphics[width=\textwidth]{figures/starcoderplus_qdep_line_plot.pdf}
    \caption{StarCoder+.}
    \label{linc-paper:fig:starcoder-qdep}
  \end{subfigure}
  \hfill
  \begin{subfigure}[b]{0.32\textwidth}
    \centering
    \includegraphics[width=\textwidth]{figures/gpt-3.5-turbo-16k-0613_qdep_line_plot.pdf}
    \caption{GPT-3.5.}
    \label{linc-paper:fig:gpt-3.5-qdep}
  \end{subfigure}
  \hfill
  \begin{subfigure}[b]{0.32\textwidth}
    \centering
    \includegraphics[width=\textwidth]{figures/gpt-4-0613_qdep_line_plot.pdf}
    \caption{GPT-4.}
    \label{linc-paper:fig:gpt-4-qdep}
  \end{subfigure}
  \caption{Accuracy per necessary proof depth in ProofWriter. Accuracies reported are for bootstrapped $10$-way majority vote, and shaded areas cover $\pm 1$ bootstrapped standard deviation. Black, dotted lines reflect the expected success rate of guessing a random label, which is 1/3 in all subsets per our experiment design.}
  \label{linc-paper:fig:results-qdep}
\end{figure}

Finally, in Figure~\ref{linc-paper:fig:results-qdep} we analyzed each model's performance on ProofWriter as a function of the necessary proof depth.
First, we found that StarCoder+'s performance was close to random chance across all three baselines, suggesting it did not have reasoning capabilities out-of-the-box.
With \texttt{LINC}, it attains almost perfect accuracy at depths $0$ and $1$, but drops as the necessary proof length increases, likely due to some formalization inaccuracies.
As StarCoder+ is a coding model, it makes sense that its formalization abilities far exceed its reasoning abilities.
For GPT-3.5, all baselines perform above chance at depth $0$, but quickly drop down for higher proof depths. In contrast,
\texttt{LINC} allows it to achieve near-perfect performance for all proof depths, showing our approach scales even for longer reasoning chains.
We saw a similar trend for GPT-4: although the baselines were significantly stronger, reasoning difficulty still remained at larger proof depths.

\begin{figure}[th]
  \centering
  \begin{subfigure}[b]{0.48\textwidth}
    \includegraphics[width=\textwidth]{figures/matrices/gpt4_neurosym_confusion.pdf}
    \caption{Confusion matrix for \texttt{LINC}.}
    \label{linc-paper:fig:gpt_neurosym_confusion}
  \end{subfigure}
  \hfill
  \begin{subfigure}[b]{0.48\textwidth}
    \includegraphics[width=\textwidth]{figures/matrices/gpt4_cot_confusion.pdf}
    \caption{Confusion matrix for \texttt{Chain-of-Thought}.}
    \label{linc-paper:fig:gpt_cot_confusion}
  \end{subfigure}
  \hfill
  \begin{subfigure}[b]{0.48\textwidth}
    \includegraphics[width=\textwidth]{figures/matrices/gpt-4_cot_neurosym.pdf}
    \caption{Comparing the consistency of \texttt{LINC} vs. \texttt{Chain-of-Thought}.}
    \label{linc-paper:fig:gpt_cot_neurosym}
  \end{subfigure}
  \hfill
  \begin{subfigure}[b]{0.48\textwidth}
    \includegraphics[width=\textwidth]{figures/matrices/gpt-4_method_similarity.pdf}
    \caption{Similarity between incorrect predictions of each method, i.e., \texttt{(A wrong == B wrong) / (A wrong or B wrong).}}
    \label{linc-paper:fig:method_similarity}
  \end{subfigure}

  \caption{Analyzing and comparing the mistakes made by GPT-4 on the FOLIO dataset.}
  \label{linc-paper:fig:confusion_plots}
\end{figure}

%%%
\subsection{Error Analysis}
\label{linc-paper:sec:error_analysis}
Having established that \texttt{LINC} can improve performance in many settings, we now move on to our final research question:
\textit{How do the failure modes of \texttt{LINC} compare to those of in-context reasoning methods}?
We focused on comparing GPT-4+\texttt{CoT} vs. GPT-4+\texttt{LINC} on FOLIO, since their overall performance was very similar (75.3\% vs. 72.5\% average accuracy). We leave an analysis of StarCoder+'s predictions on FOLIO to the original paper appendix.

\subsubsection{Qualitative Analysis} \label{linc-paper:sec:error_qualitative}
Qualitatively, we found that \texttt{LINC} and \texttt{CoT} had completely different failure modes. We give a high-level overview and abbreviated examples of each failure mode here, leaving full detailed examples to the original paper appendix.

First, we detail the failure modes for \texttt{LINC}:

\textbf{L1: FOL fails to capture implicit information not mentioned in the premises.} Often, there is obvious information not explicitly listed in the premises that is necessary to explicitly encode in FOL in order to successfully make a desired deduction. For example, in the snippet below one must encode in FOL the implicit assumption that Harry is a person (\texttt{Person(Harry)}).

\begin{itemize}
\itemsep0.1em
    \item[] \small Premise 1: \texttt{When a person reads a book, that person gains knowledge.}
    \item[] \small FOL: \parbox[t]{0.78\linewidth}{\ttfamily\small all x. all y. (Person(x) \& Reads(x, y) \& Book(y) ->\\ Gains(x, Knowledge))}
    \item[] \small Premise 2: \texttt{Harry read the book "Walden" by Henry Thoreau.}
    \item[] \small FOL: \texttt{Reads(Harry, Walden)}
    \item[] \small Conclusion (Prover9: Uncertain): \texttt{Harry gains knowledge.}
    \item[] \small FOL: \texttt{Gains(Harry, Knowledge)}
\end{itemize}


\textbf{L2: FOL fails to capture information explicitly mentioned in the premises due to the choice of representation.} Even when information is explicitly written in the premises, the choice of how the NL is represented in FOL can lead to lost information. In the example below, the fact that Heinrich was a Nazi German politician is captured by one symbol \texttt{NaziGermanPolitician}, causing the information that he was independently Nazi, German, or a politician to be lost.
As a result, \texttt{LINC} predicted \texttt{Uncertain} instead of the ground truth label \texttt{True}.

\begin{itemize}
\itemsep0.1em
    \item[] \small Premise: \texttt{Heinrich Schmidt was a Nazi German politician.}
    \item[] \small FOL: \texttt{NaziGermanPolitician (HeinrichSchmidt)}
    \item[] \small Conclusion (Prover9: Uncertain): \texttt{Heinrich Schmidt was German.}
    \item[] \small FOL: \texttt{German(HeinrichSchmidt)}
\end{itemize}

\textbf{L3: FOL contains syntax errors.}  Across all generations, we found that the FOL expressions sometimes contained syntax errors: 38\% for StarCoderPlus, 24\% for GPT-3.5, and 13\% for GPT-4. The most common error is that the same symbol is used with multiple arities. As an example, if \texttt{Summer(July4)} and \texttt{Loves(Alex, Summer)} were both present in a FOL translation, \texttt{Summer} would have a multiple arity violation.

Next, we detail three failure modes for \texttt{CoT}:

\textbf{C1: \texttt{CoT} concludes something different than it suggests.} One common mistake we often saw with the \texttt{CoT} approach was that it generated a chain of reasoning indicating uncertainty, but then asserted that the conclusion was false. An example of this is shown below, where GPT-4 suggested there was ``no evidence'' for the conclusion but then deduced it false.

\begin{lstlisting}[breaklines=true,basicstyle=\small\ttfamily]
Premises:
Some mammals have teeth.
Platypus have no teeth.
Platypus are mammals.
Humans have teeth.

Conclusion:
Platypus are reptiles.

(GPT-4 Incorrect CoT): Let's think step by step. We want to evaluate if platypus are reptiles. We know from the premises that platypus are mammals and that they have no teeth. The fact that some mammals have teeth does not contradict the fact that platypus, although mammals, do not have teeth. Moreover, the statement about humans having teeth is unrelated to the classification of platypus. There's no evidence in the premises that would suggest platypus are reptiles. Therefore, the statement is false.
\end{lstlisting}

\textbf{C2: CoT makes incorrect logical deductions.} While there were many instances where \texttt{CoT} made incorrect logical deductions, one common deduction mistake was the structure ``if $B$ is true and $A \rightarrow B$, then $A$ is true''. In the example shown below, we saw this incorrect deduction where A is ``Mike has an F1 visa'' and B is ``Mike needs to apply for CPT or OPT if he wants to work in the US''. We note that making this type of logical mistake is often called the \emph{fallacy of the converse}, and that humans are not immune to it, either.

\begin{lstlisting}[breaklines=true,basicstyle=\small\ttfamily]
Premises:
International students in US have either an F1 visa or a J1 visa.
An international student in US with an F1 visa needs to apply for CPT or OPT if the student wants to work in the US.
Mike is an international student.
Mike needs to apply for CPT if he wants to work in the US.

Conclusion (Ground Truth: Uncertain):
Mike has an F1 visa.

(GPT-4 Incorrect CoT): Let's think step by step. We want to evaluate if Mike has an F1 visa. From the premises, we know that Mike is an international student and needs to apply for CPT if he wants to work in the US. We also know that an international student in the US with an F1 visa needs to apply for CPT or OPT if the student wants to work in the US. So, it is clear that Mike has an F1 visa since he fits the criteria. Therefore, the conclusion is true.
\end{lstlisting}

\textbf{C3: \texttt{CoT} fails to find complex paths of reasoning.} We found that with all three models, \texttt{CoT} failed when the path of reasoning necessary to make the deduction was complex. Sometimes, \texttt{CoT} had difficulty getting started, and other times, it got stuck in the middle of a reasoning chain.

\subsubsection{Quantitative Analysis}

\textbf{1. Compared to \texttt{CoT}, \texttt{LINC} has worse recall but better precision on \texttt{True}/\texttt{False} predictions.} To see this,
we plotted the confusion matrices in Figure~\ref{linc-paper:fig:gpt_cot_confusion} (\texttt{CoT}) and Figure~\ref{linc-paper:fig:gpt_neurosym_confusion} (\texttt{LINC}).
Looking just at the distributions of predicted labels of the two methods, we saw that \texttt{CoT} predicted 32\% \texttt{True}, 27\% \texttt{False}, and 41\% \texttt{Uncertain}, while \texttt{LINC} predicted 24\% \texttt{True}, 17\% \texttt{False}, and 57\% \texttt{Uncertain} (with 2\% of predictions throwing an error).
Notably, we observed that \texttt{LINC} predicted \texttt{Uncertain} much more frequently than \texttt{CoT} (57\% vs. 41\%).
To understand why, note that the translation from natural language to FOL is a lossy process: recall that in L1 and L2, we saw that the information conveyed through the FOL is sometimes a subset of the information in the original premises. Removing pieces of crucial information that were on the critical path to deducing \texttt{True}/\texttt{False} may then leave an uncertain conclusion.
At the same time, while the FOL translations sometimes do not retain all of the information in the NL, they rarely contain false information that was not provided in the original premises.
Therefore, \texttt{LINC}'s precision when predicting \texttt{True} or \texttt{False} is very high (93\%) compared to that of \texttt{CoT} (81\%), but this comes at the cost of lower recall on \texttt{True}/\texttt{False} predictions (60\% for \texttt{LINC} vs. 75\% for \texttt{CoT}).

\textbf{2. \texttt{LINC} and \texttt{CoT} mispredict on different examples.} Earlier in Sec. \ref{linc-paper:sec:error_qualitative}, we saw that \texttt{LINC} and \texttt{CoT} exhibited different failure modes, which suggested they should fail on different examples. Indeed, we found that this was the case in our experiments on FOLIO:
Figure~\ref{linc-paper:fig:gpt_cot_neurosym} shows a $2\times2$ confusion matrix which compares whether or not each method's prediction was correct.
We observed that out of the $24+29+21=74$ samples where at least one method made an incorrect prediction, only 21 were shared.
On a closer examination of these 21 samples, we found that 16 were ambiguous or incorrect in their specification (details in the original paper appendix), so \textit{the two methods only agreed on 5 well-formed samples}.
This suggests that \texttt{LINC} and \texttt{CoT} are complementary methods which fail under distinct circumstances.

\textbf{3. Mispredictions of in-context reasoning baselines are more similar to each other than they are with mispredictions of \texttt{LINC}.} As an extension of the previous analysis, we next investigated the correlation between the mispredictions of each pair of methods.
To do so, we define a similarity score between two methods $A$ and $B$ as follows:
Given a dataset $D$ with $N$ rows and ground truth labels $\{A_i\}_{i=1}^N$ and $\{B_i\}_{i=1}^N$ from two methods $A$ and $B$, we define
$$\text{sim}_D(A, B) \triangleq \frac{\sum_{i=1}^N \mathbf{1}\left[ A_i = B_i \ne R_i \right] }{\sum_{i=1}^N \mathbf{1}\left[ A_i \ne R_i \text{ or } B_i \ne R_i\right]}$$
In words, $\text{sim}_D(A,B)$ measures the number of instances where $A$ and $B$ are wrong in identical ways vs. the number of instances where \emph{at least one} of them is wrong.

Figure~\ref{linc-paper:fig:method_similarity} shows the pairwise similarity between our four methods, highlighting that the similarity between \texttt{LINC}'s mispredictions and the other methods' mispredictions ($0.14$, $0.21$, $0.22$) is much lower than the similarity between any pair of the in-context reasoning methods ($0.52$, $0.54$, $0.56$).
These results suggest that for GPT-4 on FOLIO, \texttt{LINC} was the only method we evaluated which significantly altered the ways in which the model \emph{fails} to reason.


\section{Discussion}
The four examples in this chapter highlight how evaluations can go beyond ranking models and
reveal deeper insights about model behavior. In each case, analyzing the underlying behavior of
models revealed nuances hidden by aggregate scores.
Through CRUXEval, we learned that distillation may not fully lead to transfer between
different programming abilities and saw the potential and limitations of chain-of-thought
prompting and fine-tuning.
In the Counterfeit Conundrum, we saw several counterintuitive trends about the
properties of counterfeit samples and observed that counterfeits are a relatively general phenomenon and
not just a failure mode of weaker models or harder problems. 
In LiveCodeBench, we saw how post-trained models seemed to be overfit to HumanEval, suggesting 
insufficiencies in the diversity of post-training data distributions.
Finally, through LINC, simply looking at scores would show that LINC generally outperformed chain-of-thought
prompting but would miss the nuance that the two approaches were actually complementary and failed
in distinct ways. 
By analyzing the underlying successes and failures in each example, we gained a better high-level understanding
of where methodologies shine and fall short, identifying opportunities for future scientific inquiry.
